\section{Overall Structure of the Proposed Thesis}
\label{sec:structure}

The final written thesis will be divided into the following chapters:

\begin{enumerate}
    \item \textbf{Introduction.} Explains the problem, why it matters, and what this research is trying to find out. Gives the reader a clear picture of what the whole thesis is about.

    \item \textbf{Literature Review.} Summarizes what other researchers have already studied and found on related topics. Explains what is still missing from existing research and presents the three main questions this thesis will answer.

    \item \textbf{Methodology.} Describes exactly how the experiment was set up and carried out. Explains how the test documents were created, how the three local AI models and the online service were tested, how three different instruction versions were used for each system and why, how accuracy was measured using the ANLS* metric, and how the cost of each option was calculated using measured electricity data and straight-line depreciation.

    \item \textbf{Results.} Shows all the data collected from the experiments: how accurate each AI was (across all three instruction versions), how fast each one was, how much each one cost per document, and at what point each local model becomes cheaper than the online service. Charts and tables will be included to make the results easy to understand.

    \item \textbf{Discussion.} Explains what the results actually mean. Discusses whether the findings answer the research questions, what the limitations of the study are, and what future researchers could do next.

    \item \textbf{Conclusion.} A short final summary of the most important findings and what they mean for organizations thinking about using AI for document processing.

    \item \textbf{References.} A full list of all the research papers and sources used in the thesis.

    \item \textbf{Appendices.} Extra supporting material, including the exact instruction versions given to each AI during the test, sample documents from the dataset, the answer key format, the Python script used for scoring, and the full raw results tables.
\end{enumerate}
